% Generated by scripts/assemble.py; corrections belong in source/corrections.json.
\examhead{132}{华中科技大学 2024 年数学分析试卷}{125}
% p126-b12; source PDF page 126
\par\Needspace{5\baselineskip}
\blockanchor{p126-b12}
\majorquestion{一}计算题.\par
% p126-b13; source PDF page 126
\blockanchor{p126-b13}
\question{1}计算极限 $\lim_{x\to 0^{+}}\frac{1}{x}\arccos \frac{\sin x}{x}.$\par
% p126-b14; source PDF page 126
\blockanchor{p126-b14}
\question{2}计算极限 \displayformula{\lim_{n\to \infty}\left(1 + \frac{1}{n}\right)\left(1 + \frac{1}{2^3n}\right)^{2^3}\left(1 + \frac{1}{3^3n}\right)^{3^3}\dots \left(1 + \frac{1}{n^3n}\right)^{n^3}}.\par
% p126-b15; source PDF page 126
\par\Needspace{6\baselineskip}
\blockanchor{p126-b15}
\question{3}已知 $f$ 在 $\mathbb{R}$ 上具有连续导数, $L$ 是从 $A\left(3, \frac{2}{3}\right)$ 到 $B(1, 2)$ 的曲线段 $y = -\frac{5}{9} x^2 + \frac{14}{9} x + 1$, 求\par
\nopagebreak[4]
% p126-b16; source PDF page 126
\blockanchor{p126-b16}
\displayformula{I = \int_ {L} \frac {1 + y ^ {2} f (x y)}{y} \mathrm{d} x + \frac {x}{y ^ {2}} (y ^ {2} f (x y) - 1) \mathrm{d} y.}
% p126-b17; source PDF page 126
\blockanchor{p126-b17}
\question{4}计算极限 $\lim_{n\to \infty}\frac{1}{n}\int_0^{n\pi}\frac{x}{1 + n^2\cos^2x}\mathrm{d}x.$\par
% p126-b18; source PDF page 126
\par\Needspace{5\baselineskip}
\blockanchor{p126-b18}
\majorquestion{二}解答题.\par
% p126-b19; source PDF page 126
\blockanchor{p126-b19}
\question{5}用下列两种方法求 $I = \iint_{S} h(x, y, z) \, \mathrm{d}S$，其中 $S: x^2 + y^2 + z^2 = 1$，$h(x, y, z) = x^4 + y^2 z^2$.\par
% p126-b20; source PDF page 126
\blockanchor{p126-b20}
(1) 用球坐标变换.\par
% p126-b21; source PDF page 126
\par\Needspace{6\baselineskip}
\blockanchor{p126-b21}
(2) 用 Gauss 公式和齐次函数 Euler 公式\par
\nopagebreak[4]
% p126-b22; source PDF page 126
\blockanchor{p126-b22}
\displayformula{x f _ {x} + y f _ {y} + z f _ {z} = n f.}
% p126-b23; source PDF page 126
\blockanchor{p126-b23}
其中 f 满足  $f(tx,ty,tz)=t^{n}f(x,y,z)$ ，n 为正整数.\par
% p126-b24; source PDF page 126
\blockanchor{p126-b24}
\question{6}解答如下问题:\par
% p127-b00; source PDF page 127
\blockanchor{p127-b00}
(1) 设 $f$ 为 $[1, +\infty)$ 上的非负递减函数, 记 $x_{n} = \sum_{k=1}^{n} f(k) - \int_{1}^{n} f(t) \mathrm{d}t$, 证明 $\{x_{n}\}$ 收敛.\par
% p127-b01; source PDF page 127
\blockanchor{p127-b01}
(2) 计算 $\sum_{n=1}^{\infty} (-1)^{n-1} \frac{\ln n}{n}$ (可能会用到 $\sum_{k=1}^{n} \frac{1}{k} = \gamma + \ln n + o(1) (n \to \infty)$, 其中 $\gamma$ 为常数).\par
% p127-b02; source PDF page 127
\blockanchor{p127-b02}
\question{7}设正整数 $n, k$ 满足 $n \geq k$, 记 $b(n, k) = \int_{\frac{1}{4}}^{1} C_n^k t^{k-1}(1 - t)^{n-k} \mathrm{d}t$, 其中 $C_n^k = \frac{n!}{k!(n-k)!}$.\par
% p127-b03; source PDF page 127
\blockanchor{p127-b03}
(1) 证明: $\lim_{n\to \infty}\sum_{k = 1}^{n}b(n,k) = 2\ln 2.$\par
% p127-b04; source PDF page 127
\blockanchor{p127-b04}
(2) 求极限 $\lim_{n\to \infty}\frac{n4^n}{3^n}\left(2\ln 2 - \sum_{k = 1}^n b(n,k)\right)$\par
% p127-b05; source PDF page 127
\blockanchor{p127-b05}
\question{8}设  $u_{n}(x)$ ,  $S(x)$  均为  $[a, b]$  上的可积函数, 其中  $n = 1, 2, \cdots$ .\par
% p127-b06; source PDF page 127
\par\Needspace{6\baselineskip}
\blockanchor{p127-b06}
(1) 证明: 对 $S_{n}(x) = \sum_{k=1}^{n} u_{k}(x)$ 有\par
\nopagebreak[4]
% p127-b07; source PDF page 127
\blockanchor{p127-b07}
\displayformula{\lim _ {n \rightarrow \infty} \int_ {a} ^ {b} | S _ {n} (x) - S (x) | \mathrm{d} x = 0.}
% p127-b08; source PDF page 127
\blockanchor{p127-b08}
则逐项积分得到的级数 $\sum_{n=1}^{\infty} \int_{a}^{x} u_n(t) \, \mathrm{d}t$ 在 $[a, b]$ 上一致收敛于 $\int_{a}^{x} S(t) \, \mathrm{d}t$.\par
% p127-b09; source PDF page 127
\par\Needspace{6\baselineskip}
\blockanchor{p127-b09}
(2) 说明函数项级数 $\sum_{n=0}^{\infty} x^n \sin(\pi x)$ 在 $[0,1]$ 上不一致收敛, 用 (1) 证明\par
\nopagebreak[4]
% p127-b10; source PDF page 127
\blockanchor{p127-b10}
\displayformula{\sum_ {n = 0} ^ {\infty} \int_ {0} ^ {1} x ^ {n} \sin (\pi x) \mathrm{d} x = \int_ {0} ^ {1} \frac {\sin (\pi x)}{x} \mathrm{d} x.}
% p127-b11; source PDF page 127
\par\Needspace{5\baselineskip}
\blockanchor{p127-b11}
\majorquestion{三}证明题.\par
% p127-b12; source PDF page 127
\blockanchor{p127-b12}
\question{9}设 $a_0 = 1, a_{n+1} = a_n e^{\frac{1}{a_n}}, p > 2$, 证明: 级数 $\sum_{n=2}^{\infty} \frac{a_n - n}{n (\ln n)^p}$ 收敛.\par
% p127-b13; source PDF page 127
\par\Needspace{6\baselineskip}
\blockanchor{p127-b13}
\question{10}设不恒为 0 的函数 $f(x)$ 在 $[0,1]$ 上连续, 且 $\int_0^1 f(x)\mathrm{d}x = 0$. 证明: 存在 $c\in (0,1)$, 使得\par
\nopagebreak[4]
% p127-b14; source PDF page 127
\blockanchor{p127-b14}
\displayformula{c ^ {2} f (c) = \int_ {0} ^ {c} (\alpha x + x ^ {2}) f (x) \mathrm{d} x, 0 <   \alpha <   \frac {4}{3}.}
% p127-b15; source PDF page 127
\par\Needspace{6\baselineskip}
\blockanchor{p127-b15}
\question{11}设常数 $\beta \geq 1$, 证明对任意的 $x > 0$, 成立\par
\nopagebreak[4]
% p127-b16; source PDF page 127
\blockanchor{p127-b16}
\displayformula{\frac {x ^ {\beta}}{\sqrt {3}} \geq \frac {1}{x + \sqrt {3} - 1} + \frac {1 + \sqrt {3} \beta}{3} \ln x.}
% p127-b17; source PDF page 127
\blockanchor{p127-b17}
\question{12}解答如下问题:\par
% p127-b18; source PDF page 127
\blockanchor{p127-b18}
(1) 证明二元函数 $z = f(P)$ 在 $D \subset \mathbb{R}^2$ 上一致连续的充分必要条件是对 $D$ 上任意两点列 $\{P_n\}$, $\{Q_n\}$, 只要距离 $d(P_n, Q_n) \to 0 (n \to \infty)$, 就有 $f(P_n) - f(Q_n) \to 0 (n \to \infty)$.\par
% p127-b19; source PDF page 127
\par\Needspace{6\baselineskip}
\blockanchor{p127-b19}
(2) 判断函数\par
\nopagebreak[4]
% p127-b20; source PDF page 127
\blockanchor{p127-b20}
\displayformula{g (x, y) = \left\{ \begin{array}{l l} \frac {x ^ {4} y ^ {4}}{x ^ {4} + y ^ {4}}, & x ^ {4} + y ^ {4} \neq 0; \\ 0, & x ^ {4} + y ^ {4} = 0. \end{array} \right.}
% p127-b21; source PDF page 127
\blockanchor{p127-b21}
在 $\mathbb{R}^2$ 上的一致连续性.\par
